\documentclass[9pt,twocolumn,twoside]{opticajnl}
\usepackage{siunitx}
\journal{opticajournal} 
\setboolean{shortarticle}{true}

\usepackage{lineno}
\usepackage{graphicx}
\usepackage{array}      % For the >{...} syntax
\usepackage{tabularx}   % For the tabularx environment
\usepackage{multirow}   % For the \multirow command
\usepackage{booktabs}   % For \toprule, \midrule, \bottomrule

% \linenumbers % Turn off line numbering for Optica Open preprint submissions.

\usepackage{subcaption}

%\title{A Hierarchical Algorithm for Stop-Detection in Trajectory Data Without Parameter Tuning}
\title{Parameter-free robust stop-detection with HDBSCAN}

\author[1]{Andres Mondragon}
\author[1,*]{Francisco Barreras}

\affil[1]{Department of Computer and Information Science, University of Pennsylvania, 3333 Chestnut Street, Philadelphia, PA 19104}

%\affil[*]{fbarrer@sas.upenn.edu}

\setboolean{displaycopyright}{false} % Do not include copyright or licensing information in submission.

\begin{document}

% Title: A Hierarchical Algorithm for Stop-Detection in Trajectory Data Without Parameter Tuning
\maketitle
% \begin{figure}[ht]
% \vspace{-1.1em}
% \centering
% \includegraphics[width=0.8\linewidth]{figures/problem-explanation.png}
% \caption{\small Subset of a trajectory in the Garden City sandbox. Each row represents a building, and each gray rectangle a true stop. The detected stops (blue lines) exhibit the types of errors described in Table \ref{tab:stop_definitions}.}
% \label{fig:problems}
% \end{figure}
As human mobility science applications continue to mature, there are growing calls for validation of the algorithms employed in data processing \cite{DESOJO2026102397, barreras2024exciting}. In particular, large-scale GPS mobility data are known to have high temporal sparsity and spatial noise \cite{bahr2022missing, graser2021exploratory, mccool2024maximum, sanchez2026correcting}, which can hinder the performance of key algorithms such as stop detection. Furthermore, multiple algorithms exist, guidance on which is better suited to different types of data is limited \cite{safegraph_visit_attribution, birant2007st, hariharan2004project, xiang2016extracting, chen2014tdbscan, aslak2020infostop}, and how to choose the appropriate parameters remains an open question.

To make these choices more transparent, we present a set of interactive dashboards for stop detection in human mobility data. The dashboards organize commonly used methods into sequential and density-based families and describe the assumptions and parameters that govern how each method converts timestamped GPS pings into stops. Users can apply two algorithms to the same trajectory, vary parameters such as the maximum temporal gap, spatial threshold, minimum point count, and minimum dwell duration, and inspect the resulting stops side by side on a spatial map and temporal barcode. This trajectory-level view makes errors that can be difficult to identify from summary statistics immediately visible, including missed visits and stops that are split across multiple clusters or merged across nearby locations.

These local comparisons are paired with an aggregate evaluation dashboard that traces precision, recall, F1 score, and the fractions of missed, merged, and split stops across parameter sweeps. Connecting individual trajectories with summary performance allows researchers to see both where an algorithm fails and how consistently that behavior appears across a larger set of trajectories. The dashboards are intended as exploratory and educational tools rather than as a universal ranking of stop-detection methods: they show how algorithm and parameter choices propagate from raw pings to detected visits and downstream mobility measures, providing a more transparent basis for selecting and reporting a stop-detection pipeline.

\begin{table}[ht]
\vspace{-1.2em}
\centering
\caption{\small Definitions of stop‐detection terms and parameters.}
\label{tab:stop_definitions}
\small               % optional: shrink font a bit
\begin{tabularx}{\columnwidth}{@{} l X @{}}
\toprule
\textbf{Term} & \textbf{Definition} \\ \midrule
True stop                & Part of a trajectory with stationary behavior.\\
Detected stop            & Pings clustered by a stop-detection algorithm.\\
Missing stop             & A true stop not detected by the algorithm.\\
Split stop               & A true stop with separate parts correctly detected.\\
Merged stop              & A true stop partly detected as a wrong location.\\
Invalid stop set         & A group of stops with overlapping time intervals.\\
$\Delta_{\texttt{max}}$ & Maximum diameter allowed for a stop in Lachesis.\\
$\epsilon$              & Max. distance between neighbor pings in DBSCAN.\\
\texttt{max\_gap} & Max possible time gap between neighbor pings.\\
\texttt{min\_pts}        & Min neighbors for a ping to be `dense' (core).\\
\bottomrule
\end{tabularx}
\vspace{-1.25em}
\end{table}

\section{Related Work}

Stop-detection algorithms can broadly be classified into \emph{sequential} or \emph{density-based}. The former---e.g. as in 'project Lachesis' \cite{hariharan2004project}---process pings chronologically, assimilating pings into a stop until a stopping condition is triggered---such as the stop exceeding a diameter $\Delta_{\texttt{max}}$ or having a time gap larger than \texttt{max\_gap}. While simple and efficient, these methods are highly sensitive to measurement noise: a single outlier can end a stop prematurely or cause unnecessary splits.

Density-based methods, typically variations of DBSCAN \cite{ester1996density}, group pings based on spatial proximity---regardless of time---based on the reachability graph of the set of 'core' points---those with at least $\texttt{min\_pts}$ neighbors within a radius $\epsilon$. It is well-known that these methods are more robust to noise, and can produce clusters of arbitrary size. They, however, depend on choosing a good radius $\epsilon$, which is a more abstract parameter and dependent on signal density---which is heterogeneous over users and over time. Moreover, applying DBSCAN directly to GPS data neglecting the time dimension would certainly result in invalid stops that have temporal overlaps---a problem known as the 'many-visits-to-one-location' problem. Some variants address this problem through pre- and post-processing of the data, such as restricting the graph of neighbors--- based on a time threshold \texttt{max\_gap} as in \texttt{ta-dbscan} \cite{birant2007st} or dynamically T-DBSCAN \cite{chen2014tdbscan}---or simply splitting overlapping stops ex-post. Alternatively, some sequential algorithms implement 'DBSCAN-like' features to increase robustness, for instance those algorithms used by companies like SafeGraph \cite{safegraph_visit_attribution} or Cuebiq---which uses a variant of the SOC algorithm \cite{xiang2016extracting}.
Unlike DBSCAN, which requires a fixed $\epsilon$, \texttt{hdbscan} \cite{campello2013density} builds a hierarchy of clusters for all values of $\epsilon$ simultaneously. Each level of the hierarchy corresponds to a DBSCAN clustering at a particular $\epsilon$. A final clustering is selected by choosing an `antichain' in the hierarchy---a set of non-overlapping clusters from different levels. This approach effectively clusters points ``as if'' using different values of $\epsilon$ for different regions of the domain. Our method adapts \texttt{hdbscan} to spatiotemporal data by modifying the cluster selection step and adding post-processing to avoid invalid stops, as described next.
% MORE EXTENDED PAPER SHOULD ALSO CONSIDER \cite{palma2008clustering, zimmermann2009finding, tran2011robust}

\section{Data and Methods}
% Two-column, transposed, decimal-aligned version
\begin{table*}[t]
\captionsetup{justification=centering,singlelinecheck=true}
\vspace{-0.6em}
\centering
\caption{\small Performance of \texttt{hdbscan} against baselines (best parameterization), transposed so stop types become column groups.}
\label{tab:performance_by_category_wide}
\small
\renewcommand{\arraystretch}{0.95}
\setlength{\tabcolsep}{5pt}
\setlength{\extrarowheight}{0.25ex}
\sisetup{table-number-alignment=center,detect-weight=true,mode=text}

\begin{tabular}{@{}l
S[table-format=1.2]S[table-format=2.2]
S[table-format=1.2]S[table-format=2.2]
S[table-format=1.2]S[table-format=2.2]
S[table-format=1.2]S[table-format=2.2]
S[table-format=1.2]S[table-format=2.2]
S[table-format=1.2]S[table-format=2.2]@{}}
\toprule
& \multicolumn{8}{c}{\textbf{building\_type}} & \multicolumn{4}{c}{\textbf{dwell\_length}} \\
\cmidrule(lr){2-9} \cmidrule(l){10-13}
\textbf{Algorithm}
& \multicolumn{2}{c}{home}
& \multicolumn{2}{c}{park}
& \multicolumn{2}{c}{retail}
& \multicolumn{2}{c}{work}
& \multicolumn{2}{c}{high (>3h)}
& \multicolumn{2}{c}{low (<1h)} \\
\cmidrule(lr){2-3}\cmidrule(lr){4-5}\cmidrule(lr){6-7}\cmidrule(lr){8-9}\cmidrule(lr){10-11}\cmidrule(l){12-13}
& {F1} & {\scriptsize(\% Oracle)}
& {F1} & {\scriptsize(\% Oracle)}
& {F1} & {\scriptsize(\% Oracle)}
& {F1} & {\scriptsize(\% Oracle)}
& {F1} & {\scriptsize(\% Oracle)}
& {F1} & {\scriptsize(\% Oracle)} \\
\midrule
lachesis
& 0.50 & 72.97
& 0.21 & 64.83
& 0.15 & 60.14
& 0.40 & 80.50
& 0.48 & 73.91
& 0.15 & 58.72 \\
\textbf{ta-hdbscan}
& \bfseries 0.55 & \bfseries 81.28
& \bfseries 0.24 & \bfseries 79.28
& \bfseries 0.15 & \bfseries 56.66
& \bfseries 0.44 & \bfseries 88.42
& \bfseries 0.53 & \bfseries 82.46
& \bfseries 0.14 & \bfseries 56.46 \\
tadbscan
& 0.53 & 80.02
& 0.17 & 55.97
& 0.13 & 46.15
& 0.41 & 81.94
& 0.51 & 80.57
& 0.11 & 43.76 \\
\bottomrule
\end{tabular}
\vspace{0.65em}
\label{tab:performance_by_category}
\end{table*}

\begin{figure*}[hbt!]
\vspace{-1.05em}
\centering
\captionsetup{singlelinecheck=false, justification=justified}
\centering
\includegraphics[width=\linewidth]{\detokenize{figures/f1_score_as_pct_of_oracle.png}}
\caption{\small Percentage of stops with errors of each type. \texttt{lachesis\_coarse}: $(\Delta_{\texttt{max}}=40\text{m}, \texttt{max\_gap}=240\text{min})$, \texttt{lachesis\_fine}: $(\Delta_{\texttt{max}}=25\text{m}, \texttt{max\_gap}=120\text{min})$; \texttt{ta-dbscan\_coarse}: $(\epsilon=30\text{m}, \texttt{max\_gap}=240\text{min}, \texttt{min\_pts}=2)$, and \texttt{ta-dbscan\_fine}: $(\epsilon=20\text{m}, \texttt{max\_gap}=120\text{min}, \texttt{min\_pts}=3)$. The rightmost panel shows F1 score for minute-by-minute prediction as a percentage of the oracle's F1 score.}
\label{fig:errors_per_stop}
\vspace{-0.7em}
\end{figure*}


\noindent\textbf{TA-HDBSCAN.} \texttt{hdbscan} implements a greedy algorithm that selects the `antichain' of clusters with the maximum sum of \emph{cluster stabilities} \cite{campello2013density}---a metric of the 'range of persistence' of a cluster as epsilon shrinks. The 'true' modes of the density are clusters that persist over a wide range, while 'spurious' clusters disappear quickly in comparison. In \texttt{ta-hdbscan}, we adapt this procedure to incorporate time, defining (and discarding) clusters as spurious when they do not have a minimum duration, and using a definition of cluster stability that gives less weight to the stability of 'redundant' pings in a short burst, compared to more 'isolated' pings which better support cluster.

\noindent\textbf{Baseline algorithms.} We compare our method against \texttt{lachesis} \cite{hariharan2004project} and \texttt{ta-dbscan}\cite{birant2007st}, each evaluated with two parameterizations: one `coarse' (better suited for sparse signals) and one `fine' (better for dense signals).  For reference, we also implement an \texttt{oracle} baseline with perfect knowledge of which location (and thus stop) each ping is associated with. This so we can control for how mere sparsity in the data can severely affect the performance of even the best possible algorithm. Accordingly, to isolate the performance gains of one algorithm over another, The F1-scores in Table \ref{tab:performance_by_category} are expressed as a percentage of \texttt{oracle}'s F1.

\noindent\textbf{Post-processing.} \texttt{ta-dbscan} and \texttt{ta-hdbscan} apply a simple recursive post-processing to break-apart invalid stops: for any set of overlapping stops, we find a minimal segment in the middle of the trajectory such that the 'tails' do not have invalid stops, and re-apply the algorithm on that middle and on each tail. Recursing if the middle also has invalid stops. 

\noindent\textbf{Error metrics.} The problems of ``merging'', ``splitting'' and ``missing'' (Table \ref{tab:stop_definitions}) are measured by counting true stops that (temporally) intersect: a stop attributed to another location; multiple stops attributed to the correct location; or no detected stop. We also frame stop detection as a multi-class prediction task, assigning each minute of a trajectory to a stop (or gap) then to a building ID (or 'street'). This allows computation of per-minute precision, recall, and F1-scores. While F1-scores for \texttt{ta-hdbscan} are overall higher, we present F1-scores (as percentage of oracle) for different parts of the trajectory, focusing on location types with distinctive sizes (e.g., parks are bigger than homes) and average dwell times (e.g., visits to workplaces are longer than to retail). We also analyze performance for short vs long durations of a true stop. A variable that can be directly tied to differences in algorithm performance is the average ``density'' of the trajectory, defined as the inverse of the distance to the closest neighbor averaged over all pings; thus we can see whether \texttt{ta-hdbscan} outperforms baselines in different density regimes. This corresponds to the ``core-density'' attribute used in \texttt{ta-hdbscan}.


\section{Results and Conclusion}

Figure~\ref{fig:errors_per_stop} shows that \texttt{ta-hdbscan} achieves the lowest missed-stop fraction and highest F1 score among all algorithms, recovering more true locations without the need for parameter tuning. As expected, sequential methods like \texttt{lachesis} show more splitting (due to sensitivity to noise), while density-based methods tend to merge stops. \texttt{ta-hdbscan} avoids this trade-off and performs comparably to the \texttt{oracle} baseline.

Table~\ref{tab:performance_by_category} breaks down F1 scores by stop type. \texttt{ta-hdbscan} consistently outperforms other methods in homes, parks, and workplaces, and in long-dwell stops, all of which contribute substantially to mobility metrics. The exception is short stops and visits to retail locations---also characterized by short dwell times. Interestingly, \texttt{lachesis} outperforms \texttt{ta-dbscan}  in longer stops and buildings with large footprint (like 'park'), showing that perhaps the density constraints come at a cost in regimes in which noise is less of a problem.

\noindent\textbf{Conclusions.} We propose a parameter-free stop-detection algorithm based on \texttt{hdbscan} that overcomes the limitations of parameter setting and lack of robustness. The results are limited by the use of synthetic data and its possible idiosyncrasies, but they represent a proof of concept and possible improvement over existing processing pipelines. 

\begin{backmatter}
\bmsection{Funding} This work was possible with the support of the National Science Foundation, NSF award 2341932.

% \bmsection{Acknowledgment} TODO

%\bmsection{Disclosures} The authors declare no conflicts of interest.
\bmsection{Data availability}  Our dataset consists of 3 week-long trajectories from 200 users generated using the Garden City trajectory generator \cite{li2024garden}. 'Macro' trajectories follow an Exploration and Preferential Return model, and 'Micro' trajectories follow a constrained brownian motion. Available at \url{https://github.com/Watts-Lab/nomad/tree/main/examples/research}.
\end{backmatter}

% \begin{figure*}[ht]
% \caption{F1 score for each algorithm compared against oracle. The scores are presented as a scatter plot with "trajectory completeness" on the x-axis.}
% \label{fig:problems}
% \centering
% \includegraphics[width=\linewidth]{figures/f1_vs_q.png}
% \end{figure*}

% Bibliography
\bibliography{refs}


\end{document}
